\documentclass[10pt,a4paper]{amsart}
\usepackage{amsmath,amssymb,amsthm,graphicx}
\usepackage{verbatim}
\newenvironment{reference}{\comment}{\endcomment}
%\newenvironment{reference}{}{}
\newenvironment{slogan}{\comment}{\endcomment}
\newenvironment{history}{\comment}{\endcomment}

% For commutative diagrams we use Xy-pic
\usepackage[all]{xy}

% We use 2cell for 2-commutative diagrams.
\xyoption{2cell}
% Theorem environments.
%
\theoremstyle{plain}
\newtheorem{theorem}[subsection]{Theorem}
\newtheorem{proposition}[subsection]{Proposition}
\newtheorem{lemma}[subsection]{Lemma}

\theoremstyle{definition}
\newtheorem{definition}[subsection]{Definition}
\newtheorem{example}[subsection]{Example}
\newtheorem{exercise}[subsection]{Exercise}
\newtheorem{situation}[subsection]{Situation}

\theoremstyle{remark}
\newtheorem{remark}[subsection]{Remark}
\newtheorem{remarks}[subsection]{Remarks}

\numberwithin{equation}{subsection}

% Macros
%
\def\lim{\mathop{\mathrm{lim}}\nolimits}
\def\colim{\mathop{\mathrm{colim}}\nolimits}
\def\Spec{\mathop{\mathrm{Spec}}}
\def\Hom{\mathop{\mathrm{Hom}}\nolimits}
\def\Ext{\mathop{\mathrm{Ext}}\nolimits}
\def\SheafHom{\mathop{\mathcal{H}\!\mathit{om}}\nolimits}
\def\SheafExt{\mathop{\mathcal{E}\!\mathit{xt}}\nolimits}
\def\Sch{\mathit{Sch}}
\def\Mor{\mathop{\mathrm{Mor}}\nolimits}
\def\Ob{\mathop{\mathrm{Ob}}\nolimits}
\def\Sh{\mathop{\mathit{Sh}}\nolimits}
\def\NL{\mathop{N\!L}\nolimits}
\def\CH{\mathop{\mathrm{CH}}\nolimits}
\def\proetale{{pro\text{-}\acute{e}tale}}
\def\etale{{\acute{e}tale}}
\def\QCoh{\mathit{QCoh}}
\def\Ker{\mathop{\mathrm{Ker}}}
\def\Im{\mathop{\mathrm{Im}}}
\def\Coker{\mathop{\mathrm{Coker}}}
\def\Coim{\mathop{\mathrm{Coim}}}

% Boxtimes
%
\DeclareMathSymbol{\boxtimes}{\mathbin}{AMSa}{"02}

%
% Macros for moduli stacks/spaces
%
\def\QCohstack{\mathcal{QC}\!\mathit{oh}}
\def\Cohstack{\mathcal{C}\!\mathit{oh}}
\def\Spacesstack{\mathcal{S}\!\mathit{paces}}
\def\Quotfunctor{\mathrm{Quot}}
\def\Hilbfunctor{\mathrm{Hilb}}
\def\Curvesstack{\mathcal{C}\!\mathit{urves}}
\def\Polarizedstack{\mathcal{P}\!\mathit{olarized}}
\def\Complexesstack{\mathcal{C}\!\mathit{omplexes}}
% \Pic is the operator that assigns to X its picard group, usage \Pic(X)
% \Picardstack_{X/B} denotes the Picard stack of X over B
% \Picardfunctor_{X/B} denotes the Picard functor of X over B
\def\Pic{\mathop{\mathrm{Pic}}\nolimits}
\def\Picardstack{\mathcal{P}\!\mathit{ic}}
\def\Picardfunctor{\mathrm{Pic}}
\def\Deformationcategory{\mathcal{D}\!\mathit{ef}}
\usepackage{cleveref}
\crefname{originallemma}{original lemma}{original lemmas}
\Crefname{originallemma}{Original lemma}{Original lemmas}
\makeatletter
\newlabel{lemma-precompose-property-local-source}{{06EQ}{source}{}{sourcepointer-1}{}}
\newlabel{lemma-precompose-property-local-source@cref}{{[originallemma][0][]06EQ}{[1][0][]source}{}{}{}}
\makeatother
\crefname{originallemma}{original lemma}{original lemmas}
\Crefname{originallemma}{Original lemma}{Original lemmas}
\makeatletter
\newlabel{lemma-pullback-property-local-target}{{06EM}{source}{}{sourcepointer-2}{}}
\newlabel{lemma-pullback-property-local-target@cref}{{[originallemma][0][]06EM}{[1][0][]source}{}{}{}}
\makeatother
\crefname{originaldefinition}{original definition}{original definitions}
\Crefname{originaldefinition}{Original definition}{Original definitions}
\makeatletter
\newlabel{definition-local-source-target}{{06F8}{source}{}{sourcepointer-3}{}}
\newlabel{definition-local-source-target@cref}{{[originaldefinition][0][]06F8}{[1][0][]source}{}{}{}}
\makeatother
\title{A property smooth local on source and target, stable under postcomposition with smooth morphisms, is smooth local on the source-and-target}
\author{The Stacks Project Authors}
\date{stacks-project:spaces-descent.tex:local-snapshot-20261012; GNU FDL 1.2 or later}
\begin{document}
\maketitle
\noindent\textit{Editorial note.} All delivered mathematical text is the author's own, copied verbatim from the pinned Stacks Project chapter file spaces-descent.tex; no statement or proof step was rewritten, reordered or removed.\par\medskip
\noindent\textit{Editorial note.} The counted span cites two lemmas and one definition of the same chapter, recorded in assembly.external\_references with canonical tags as navigation pointers.\par\medskip
\noindent\textit{Allowed standard prerequisites (editorial).}\begin{itemize}
\item Smooth local on the source: precomposition stability. Stacks Project, Descent on Spaces, Lemma lemma-precompose-property-local-source. Supplies that P is stable under precomposing with smooth morphisms.
\item Smooth local on the target: base change stability. Stacks Project, Descent on Spaces, Lemma lemma-pullback-property-local-target. Supplies that P is stable under smooth base change.
\item Smooth local on the source-and-target: definition. Stacks Project, Descent on Spaces, Definition definition-local-source-target. Supplies part (3)(b) being assumed and part (3) being proved.
\end{itemize}\medskip
\ifdefined\hypertarget\hypertarget{sourcepointer-1}{}\fi
\noindent\textit{Source reference.} Original lemma 06EQ: sources/stacks-chow/tags.txt, source lines 8224--8224. Cross-chapter standard prerequisite cited by the author inside the counted span; it is used as a pointer and its own proof is not part of this item.\par\medskip
\ifdefined\hypertarget\hypertarget{sourcepointer-2}{}\fi
\noindent\textit{Source reference.} Original lemma 06EM: sources/stacks-chow/tags.txt, source lines 8221--8221. Cross-chapter standard prerequisite cited by the author inside the counted span; it is used as a pointer and its own proof is not part of this item.\par\medskip
\ifdefined\hypertarget\hypertarget{sourcepointer-3}{}\fi
\noindent\textit{Source reference.} Original definition 06F8: sources/stacks-chow/tags.txt, source lines 8242--8242. Cross-chapter standard prerequisite cited by the author inside the counted span; it is used as a pointer and its own proof is not part of this item.\par\medskip
\begin{lemma}
\label{lemma-smooth-local-source-target}
Let $S$ be a scheme.
Let $\mathcal{P}$ be a property of morphisms of algebraic spaces over $S$.
Assume
\begin{enumerate}
\item $\mathcal{P}$ is smooth local on the source,
\item $\mathcal{P}$ is smooth local on the target, and
\item $\mathcal{P}$ is stable under postcomposing with smooth morphisms:
if $f : X \to Y$ has $\mathcal{P}$ and $Y \to Z$ is a smooth morphism
then $X \to Z$ has $\mathcal{P}$.
\end{enumerate}
Then $\mathcal{P}$ is smooth local on the source-and-target.
\end{lemma}

\begin{proof}
Let $\mathcal{P}$ be a property of morphisms of algebraic spaces which
satisfies conditions (1), (2) and (3) of the lemma. By
Lemma \ref{lemma-precompose-property-local-source}
we see that $\mathcal{P}$ is stable under precomposing with
smooth morphisms. By
Lemma \ref{lemma-pullback-property-local-target}
we see that $\mathcal{P}$ is stable under smooth base change.
Hence it suffices to prove part (3) of
Definition \ref{definition-local-source-target}
holds.

\medskip\noindent
More precisely, suppose that $f : X \to Y$ is a morphism
of algebraic spaces over $S$ which satisfies
Definition \ref{definition-local-source-target} part (3)(b).
In other words, for every $x \in X$ there exists a smooth
morphism $a_x : U_x \to X$, a point $u_x \in |U_x|$ mapping to $x$,
a smooth morphism $b_x : V_x \to Y$, and a morphism $h_x : U_x \to V_x$
such that $f \circ a_x = b_x \circ h_x$ and $h_x$ has $\mathcal{P}$.
The proof of the lemma is complete once we show that $f$ has $\mathcal{P}$.
Set $U = \coprod U_x$, $a = \coprod a_x$, $V = \coprod V_x$,
$b = \coprod b_x$, and $h = \coprod h_x$. We obtain a
commutative diagram
$$
\xymatrix{
U \ar[d]_a \ar[r]_h & V \ar[d]^b \\
X \ar[r]^f & Y
}
$$
with $a$, $b$ smooth, $a$ surjective. Note that $h$ has $\mathcal{P}$
as each $h_x$ does and $\mathcal{P}$ is smooth local on the target.
Because $a$ is surjective and $\mathcal{P}$ is smooth local on the source,
it suffices to prove that $b \circ h$ has $\mathcal{P}$.
This follows as we assumed that $\mathcal{P}$ is stable under
postcomposing with a smooth morphism and as $b$ is smooth.
\end{proof}

\end{document}
